\documentclass[10pt,a4paper]{amsart}
\usepackage[margin=19mm]{geometry}
\usepackage{amsmath,amssymb,mathtools}
\usepackage[scr=rsfs,cal=euler]{mathalfa}
\usepackage{xfrac}
\usepackage[unicode,hidelinks,bookmarks=false]{hyperref}
\newcommand{\ie}{i.e.\ }
\newcommand{\cf}{cf.\ }
\newcommand{\resp}{respectively\ }
\newcommand{\Subsection}{\subsection}
\providecommand{\nobreakdash}{\nobreak\hbox{-}}
\setlength{\emergencystretch}{2em}
\multlinegap0pt
\usepackage{amssymb}
\usepackage{enumerate}
\usepackage{float}
\usepackage{xfrac}
\usepackage{mathtools}
\usepackage[shortlabels]{enumitem}
\usepackage{tikz}
\newtheorem{theorem}{Theorem}
\newtheorem{corollary}{Corollary}
\newcommand{\Ecos}{E^{\cos}}
\newcommand{\Esin}{E^{\sin}}
\newcommand{\N}{\mathbb N}
\newcommand{\R}{\mathbb R}
\newcommand{\Rp}{{\mathbb R}_{+}}
\newcommand{\Wp}{{\mathcal W}_{+}}
\newcommand{\Z}{\mathbb Z}
\renewcommand{\d}{\,\mathrm{d}}
\newcommand{\dpl}{\delta_{+}}
\newcommand{\e}{\mathrm e}
\newcommand{\fcos}{{\hat f}_{\cos}}
\newcommand{\fsin}{{\hat f}_{\sin}}
\renewcommand{\i}{\mathrm i}
\newcommand{\sign}{\mathrm{sign}}
\title{How well do fast discrete trigonometric transforms\\ work for parametric numerical integration?}
\author{Gerlind Plonka, Daniel Potts, Manfred Tasche}
\date{Source: arXiv:2609.15577v1 (CC BY 4.0)}
\begin{document}
\maketitle

\noindent\emph{Source and licence.} How well do fast discrete trigonometric transforms\\ work for parametric numerical integration?. Version arXiv:2609.15577v1 (CC BY 4.0), distributed by arXiv under the Creative Commons Attribution 4.0 International licence, http://creativecommons.org/licenses/by/4.0/, as recorded in the arXiv licence metadata for this exact version and shown on the abstract page https://arxiv.org/abs/2609.15577v1. Full licence notice: \texttt{licenses/arxiv-2609.15577-CC-BY-4.0.txt}.\par\smallskip

\noindent\textit{Editorial scope.} Counted unit: the article's theorem Ex:sech (source0.tex lines 808--817). The article's complete original proof of it is delivered with it (source0.tex lines 818--871, one \texttt{proof} environment of 54 lines). No mathematics is written, repaired or re-derived: the statement and the proof are the article's own lines, in the article's own notation, and no retained line is substituted or re-flowed.  Retained uncounted as the source-local context the proof needs: lines 189--194; lines 351--354; lines 471--492; lines 587--601; lines 680--682.  Nothing was omitted from any retained span, so the package carries no omission record.  No retained line contains a citation, so the wrapper carries no bibliography and no bibliography entry.  No retained line contains a figure environment, an image-inclusion command or drawing code, so no image is delivered and the package declares no asset dependency.  Editorial formatting only, in addition: an AMS \texttt{amsart} preamble that restates the article's own declarations and macros used by the retained lines, namely the environments \texttt{theorem}, \texttt{corollary} on separate counters, together with the standard packages those lines need.  The retained environments print in this document as Theorem 1, Corollary 1 in order of appearance, whereas the article prints the same items with the numbering of its own full document.  The complete native source of the article as distributed in the arXiv e-print for this exact version is bundled unchanged as \texttt{sources/210429/source0.tex}.
\begin{align}\label{eq:approx}
\textstyle c_{P,L} (v) \!:= \!  \tfrac{f(0)}{2P} + \tfrac{1}{P} \sum\limits_{n=1}^{LP/2} f(\frac{n}{P})  \cos(\tfrac{2\pi n v}{P}),
\quad
s_{P,L} (v) \!:=\!  \tfrac{1}{P} \sum\limits_{n=1}^{LP/2} f(\tfrac{n}{P})  \sin(\tfrac{2\pi n v}{P}),
\; v \in \Rp,
\end{align}

\begin{equation}
\label{eq:decayrate+}
\textstyle \dpl(f,L) := \max\limits_{|x| \leq L/2} \sum\limits_{k = 1}^{\infty} |f(x + k L)|\,
\end{equation}

\begin{theorem}
\label{Thm:samplingcospolynomial}
Let $L$, $P \in 2 \N$ be fixed and let $f \in \Wp$ with $\fcos \in \Wp$ be given.
Then $\fcos$ can be uniformly approximated on $[0,\,P/2]$ by the $P$-periodic
sampling cosine polynomial $c_{P,L}$ in  \eqref{eq:approx}
and  the error
estimate
\begin{equation}
\label{eq:cosmain}
 \textstyle \Ecos_{P,L}(f) =  \textstyle \frac{1}{\sqrt{L}} \max\limits_{v \in [0,\frac{P}{2}]} |\hat{f}_{\cos} (v)\! - \!c_{P,L}(v)|\leq \frac{2}{\sqrt L}\,\dpl\big(\fcos, P\big) + \sqrt L \,\dpl(f,L)
\end{equation}
holds with the (one-sided) decay rates defined in \eqref{eq:decayrate+} for $f$ and $\fcos$.

Furthermore, for $f \in \Wp$ with $\fsin \in \Wp$ and $f(0)=0$  we similarly  have
\begin{equation}
\label{eq:sinmain}
 \textstyle \Esin_{P,L}(f)= \textstyle \frac{1}{\sqrt{L}} \max\limits_{v \in [0,\frac{P}{2}]} |\hat{f}_{\sin} (v)\! - \!s_{P,L}(v)|
  \leq \frac{2}{\sqrt L}\,\dpl\big(\fsin, P\big) + \sqrt L \,\dpl(f,L).
\end{equation}
The additional assumption $f(0)=0$ in the sine case is needed for the continuity of the odd
extension of $f$ on $\R$.
\end{theorem}

\begin{corollary}\label{corlower}
Let $L$, $P \in 2 \N$ be fixed and let $f \in \Wp$ with $\fcos \in \Wp$ be given.\\
If $f$ satisfies $\sign(\fcos(\frac{P}{2})) = \sign\big(\sum_{k=1}^\infty \fcos\big(\frac{P}{2} +kP\big)\big)$, then
\begin{equation}
\label{eq:cosbelow1}
 \textstyle    \frac{1}{\sqrt{L}} \big|\fcos(\frac{P}{2}) \big| - \sqrt{L} \,\dpl(f,L)
 \leq \Ecos_{P,L}(f).
\end{equation}
If $f$ is non-negative and monotonically decreasing then
\begin{equation}
\label{eq:cosbelow2}
 \textstyle   \frac{1}{\sqrt{L}}  \int\limits_{\frac{L}{2} + \frac{1}{P}}^{\infty} f(x) \, \d x-  \frac{2}{\sqrt L}\,\dpl\big(\fcos, P\big)
 \leq \Ecos_{P,L}(f).
\end{equation}
\end{corollary}

\begin{align}\label{s0}
	s_0 := \min\limits_{j=1,\ldots,J} |\Im ( z_j)| > 0\,, \qquad k_{\max} := \max\limits_{j=1,\ldots,J} k_j\,.
\end{align}

\begin{theorem}
	\label{Ex:sech} Let $P \in 2 \N$  and $P \ge 8$. Then we find for $f(x) := \operatorname{sech}(\pi x)$ and $L=P$ the estimate
	\begin{equation}
		\label{eq:sechlower}
		\textstyle
		 \frac{1}{4\sqrt P}\,\e^{-\pi P/2} \;\leq\;
		\Ecos_{P,P}(f) \;\leq\; \frac{1.64}{\sqrt{P}}\,\e^{-\pi P/2}  .
	\end{equation}

\end{theorem}	

\begin{proof}
	The function $f(x) := \operatorname{sech}(\pi x) = \frac{2}{e^{\pi x} + e^{-\pi x}}$ is meromorphic with the
	simple poles $\i\,(k+\tfrac12)$, $k \in \Z$, and similarly as in (\ref{s0}) we set $s_0:=  \min\limits_{k \in {\mathbb Z}} |\Im( \i\,(k+\tfrac12))| = \frac{1}{2}$. The  Fourier cosine transform is given by
	$$
	\fcos(v) =  \textstyle \int\limits_0^{\infty} \operatorname{sech}(\pi x)\,\cos(2\pi x v)\,\d x
	= \tfrac12\,\operatorname{sech}(\pi v)\,, \quad v \in \Rp\,.
	$$
As in the proofs of Theorem \ref{Thm:samplingcospolynomial} and Corollary \ref{corlower},
$\sqrt{P}\,\Ecos_{P,P}(f) = \max_{v \in [0, P/2]} |\sigma_2(v)-\sigma_1(v)|$, so that
\begin{align}\label{eq:ung}
 \textstyle
\big|\sigma_1\big(\frac{P}{2}\big)\big| - \big|\sigma_2\big(\frac{P}{2}\big)\big| \le \sqrt{P} \, \Ecos_{P,P}(f)
\le   \max\limits_{v \in [0,\, P/2]} \big(\big|\sigma_1\big(v)\big| + \big|\sigma_2\big(v)\big| \big).
\end{align}
We observe that  $ \operatorname{sech} (t)
 \in [ \e^{-t} , \,  2 \e^{-t}]$ is monotonically decreasing
 for $t \ge 0$, such that
$$
 \textstyle \big|\sigma_1\big(\tfrac P2\big)\big|= \fcos\big(\tfrac P2\big) + 2 \sum\limits_{k=1}^\infty \fcos(\frac{P}{2} + kP) \ge  \tfrac12\operatorname{sech}\tfrac{\pi P}{2} \;\geq\; \tfrac12\,\e^{-\pi P/2}\,,
$$
Since $\operatorname{sech}''(t) = \operatorname{sech}(t)\big(1-2\operatorname{sech}^2 t\big)>0$
for $t > \operatorname{arsech}\frac{1}{\sqrt2} = 0.8814$, the function $\operatorname{sech}$ is
convex on $[0.89,\infty)$. For $P \ge 8$, $k \in \N$ and $v \in [0,\frac P2]$ all arguments
$\pi(kP\pm v)$ lie in $[ 4 \pi,\infty)$, so that
$v \mapsto \fcos(v+kP)+\fcos(kP-v)$ is convex
 and increasing in $v$ on $[0,\frac P2]$, hence maximal at
$v=\frac P2$. Summing over $k$ gives $|\sigma_1(v)| \le \sigma_1\big(\frac P2\big)$ and
\begin{align*}
 \textstyle \big|\sigma_1(v)\big| &\le   \textstyle
 \tfrac{1}{2}\operatorname{sech}\tfrac{\pi P}{2} +  \sum\limits_{k=1}^\infty \operatorname{sech}\big(\pi(\tfrac P2 + kP)\big)
 \le  \e^{-P\pi/2} + 2 \sum\limits_{k=1}^\infty \e^{-\frac{P\pi}{2}(1+2k)}
 < \e^{-P\pi/2} \big(1 + 10^{-10}\big)\,.
 \end{align*}
On the other hand, $\sigma_2\big(\tfrac P2\big) = \frac1P\sum_{n>P^2/2}(-1)^n\operatorname{sech}\tfrac{\pi n}{P}$
is an alternating series with monotonically decreasing terms, whence
$\big|\sigma_2\big(\tfrac P2\big)\big| \le \tfrac1P\operatorname{sech}\tfrac{\pi P}{2} \le \tfrac2P\e^{-\pi P/2}$.
 To obtain an upper bound for $\big|\sigma_2\big(v\big)\big|$ for all $v \in [0,\frac P2]$, we use
$\operatorname{sech}(t) \le 2\e^{-t}$ and estimate for every $v \in [0,\frac P2]$,
$$
\textstyle |\sigma_2(v)| \le \frac2P \sum\limits_{n>P^2/2}\e^{-\pi n/P}
= \frac{2}{P}\,\frac{\e^{-\pi P/2}}{\e^{\pi/P}-1}
= \Big( \frac{2}{\pi}\,\frac{t}{\e^{t}-1}\Big|_{t = \pi/P} \Big) \, \e^{-\pi P/2}
\le \frac{2}{\pi}\,\e^{-\pi P/2}\,,
$$
since $t \mapsto \frac{t}{\e^t-1}$ is decreasing with limit $1$ at $t=0$.
Thus we obtain from (\ref{eq:ung}) the estimate
\begin{align*}
 \textstyle
\frac{1}{\sqrt{P}} \big( \frac{1}{2} \e^{-\pi P/2}  -  \frac{2}{P} \e^{-\pi P/2} \big) \le   \Ecos_{P,P}(f)
\le   \frac{1}{\sqrt{P}}  \big( \tfrac{2}{\pi} \e^{-\pi P/2} + \e^{-\pi P/2} (1 + 10^{-10}) \big),
\end{align*}
and the assertion follows for $P\ge 8$, since $\frac12 - \frac2P \ge \frac14$ and
$\frac2\pi + 1 + 10^{-10} < 1.64$.
\end{proof}

\end{document}
